\chapter{Diseño técnico y arquitectura implementada}
\label{cap:diseno-tecnico}

\section{Vista de arquitectura}
VeriVoto implementa una aplicación web Angular/Ionic y cinco servicios
Spring: \texttt{api-gateway}, \texttt{identity-service},
\texttt{election-service}, \texttt{voter-access-service} y
\texttt{ballot-box-service}. Keycloak presta la autenticación y PostgreSQL
almacena los datos. La elección de cinco servicios y la exclusión del revoto
corresponden a D-05 y D-03; ambas fueron confirmadas por el tutor el
30/10/2026 mediante correo electrónico.

La figura~\ref{fig:arquitectura-implementada} muestra los componentes que se
ejecutan en el entorno local evaluado. A diferencia de la arquitectura
objetivo de la figura~\ref{fig:arquitectura-componentes}, no hay CDN, WAF,
KMS, cola de mensajes ni \emph{workers}: la propagación entre servicios se
resuelve con REST y una tabla \emph{outbox}. Para no recargar el diagrama se
omiten dos relaciones comunes a todos los servicios: la validación de los
tokens emitidos por Keycloak y el envío de las acciones administrativas a
\texttt{identity-service}, que agrega la trazabilidad.

\begin{figure}[htbp]
  \centering
  \resizebox{0.98\textwidth}{!}{% Arquitectura implementada en el MVP (entorno local). Misma notación que la figura de arquitectura objetivo.
\begin{tikzpicture}[
  font=\sffamily\scriptsize,
  client/.style={draw=gray!65,rounded corners=2pt,fill=gray!8,align=center,
    minimum height=.72cm,text width=6.2cm,inner sep=3pt},
  service/.style={draw=blue!50!black,rounded corners=2pt,fill=blue!5,align=center,
    minimum height=1.25cm,text width=2.55cm,inner sep=3pt},
  store/.style={draw=teal!55!black,rounded corners=2pt,fill=teal!8,align=center,
    minimum height=.7cm,text width=2.3cm,inner sep=3pt},
  external/.style={draw=gray!65,rounded corners=2pt,fill=gray!8,align=center,
    minimum height=.72cm,text width=2.75cm,inner sep=3pt},
  sync/.style={-{Stealth},thick,draw=gray!70},
  async/.style={-{Stealth},thick,dashed,draw=gray!70},
  lbl/.style={font=\sffamily\tiny,fill=white,inner sep=1pt},
  zone/.style={anchor=west,fill=white,font=\sffamily\scriptsize\bfseries,inner sep=2pt}
]
% Cliente y componentes de apoyo
\node[client] (web) at (0,0) {Aplicación web Angular/Ionic (PWA)\\flujo del votante y consola de gestión};
\node[external] (keycloak) at (-5.4,-1.5) {Keycloak\\identidad de gestión y SSO institucional};
\node[service,minimum height=.82cm,text width=3.2cm] (gateway) at (0,-1.5) {api-gateway\\punto de entrada de la API};
\node[external] (translate) at (5.4,0) {LibreTranslate (local)\\traducción de textos};
\node[external] (mail) at (5.4,-1.5) {Mailpit (SMTP local)\\enlaces y códigos};

% Servicios propios
\node[service] (identity) at (-5.4,-3.7) {identity-service\\organizaciones, roles, claves de API, trazabilidad};
\node[service] (election) at (-1.8,-3.7) {election-service\\censos, procesos, papeletas, ciclo de vida};
\node[service] (access) at (1.8,-3.7) {voter-access-service\\acceso del votante, participación, avisos};
\node[service] (ballot) at (5.4,-3.7) {ballot-box-service\\votos, recibos, cadena, escrutinio};

% Persistencia
\draw[dashed,rounded corners=4pt,draw=gray!55] (-6.95,-6.95) rectangle (3.35,-5.45);
\node[zone] at (-6.8,-6.95) {pg-main (PostgreSQL; aloja también la base de Keycloak)};
\node[store] (identitydb) at (-5.4,-6.1) {identity\_db};
\node[store] (electiondb) at (-1.8,-6.1) {election\_db};
\node[store] (accessdb) at (1.8,-6.1) {voter\_access\_db};
\draw[dashed,rounded corners=4pt,draw=gray!55] (3.85,-6.95) rectangle (6.95,-5.45);
\node[zone] at (4.0,-6.95) {pg-ballot (PostgreSQL)};
\node[store] (ballotdb) at (5.4,-6.1) {ballot\_box\_db};

% Entrada
\draw[sync] (web) -- (gateway);
\draw[sync] (web.west) -| node[lbl,pos=.25] {OIDC} (keycloak.north);
\draw[async] (web) -- node[lbl] {consola} (translate);
\draw[sync] (gateway.south) -- ++(0,-.5) -| (identity.60);
\draw[sync] (gateway.south) -- ++(0,-.5) -| (election.north);
\draw[sync] (gateway.south) -- ++(0,-.5) -| (access.120);
\draw[sync] (gateway.south) -- ++(0,-.5) -| (ballot.north);

% Entre servicios
\draw[sync] (identity.125) -- node[lbl,anchor=east,xshift=-2pt] {API de admin.} (identity.125 |- keycloak.south);
\draw[async] (election) -- node[lbl] {outbox} (access);
\draw[async] (election.south east) to[bend right=22] node[lbl,pos=.3] {outbox} (ballot.south west);
\draw[sync] (ballot) -- node[lbl] {token} (access);
\draw[sync] (access.50) |- node[lbl,pos=.75] {SMTP} (mail.west);

% Persistencia
\draw[sync] (identity) -- (identitydb);
\draw[sync] (election) -- (electiondb);
\draw[sync] (access) -- (accessdb);
\draw[sync] (ballot) -- (ballotdb);
\end{tikzpicture}
}
  \caption{Arquitectura implementada en el MVP (entorno local). Las flechas discontinuas indican propagación mediante \emph{outbox} o un componente de apoyo de la consola.}
  \label{fig:arquitectura-implementada}
\end{figure}

La separación lógica de responsabilidades es la siguiente:
\begin{description}
 \item[Gateway.] Punto de entrada de la API y encaminamiento a los servicios.
 \item[Identity.] Organizaciones, roles, claves de API y agregación de
 trazabilidad administrativa.
 \item[Election.] Censos, procesos, papeletas, ciclo de vida y propagación de
 instantáneas y estados.
 \item[Voter access.] Credenciales del votante, participación, tokens de
 emisión y resultados, notificaciones y recordatorios.
 \item[Ballot box.] Recepción de votos, recibos, cadena de auditoría,
 escrutinio, resultados y actas.
\end{description}

\section{Persistencia y separación de información}
\texttt{pg-main} contiene bases lógicas de identity, election, voter-access y
Keycloak. \texttt{pg-ballot} es una instancia separada para la urna. En
voter-access se registra la participación por día, mientras que la urna
almacena votos y eventos sin identidad ni marca temporal por voto. Esta
separación reduce la relación directa entre persona y selección; no protege
frente a un operador que controle y correlacione toda la infraestructura.

\section{Comunicación y patrón outbox}
Las llamadas entre servicios usan REST y credenciales de servicio. Las rutas
internas autorizan el cliente esperado mediante el claim \texttt{azp} (D-43).
Election persiste eventos de propagación en una tabla \emph{outbox}; un
despachador aplica reintentos e idempotencia en los receptores. El MVP no
incorpora un broker de mensajes (D-06).

\section{Autenticación y aislamiento multi-tenant}
Keycloak autentica a los usuarios de gestión. Para votantes, cada proceso
selecciona \texttt{EMAIL\_CODE} o \texttt{INSTITUTIONAL\_SSO} (D-48). En el
segundo flujo se usa Authorization Code con PKCE, \texttt{state} y
\texttt{nonce}; el correo verificado debe coincidir con la entrada del censo.
El enlace personal sigue siendo necesario. Las consultas de negocio se
restringen por \texttt{organization\_id}. El rol \texttt{INTEGRATION} ofrece
acceso de solo lectura.

\section{Decisiones arquitectónicas}
Los ADR-0001--0009 documentan, entre otros aspectos, la arquitectura
hexagonal, REST y outbox, la separación física de la urna, la generación de
recibos y la cadena de auditoría. La consolidación en cinco servicios (D-05)
y la exclusión del revoto del MVP (D-03), que queda como trabajo futuro (P9),
fueron confirmadas por el tutor el 30/10/2026 mediante correo
electrónico.

% [Bloque D] Arquitectura hexagonal y experiencia previa.
La adopción de la arquitectura hexagonal se apoya en su uso previo en
microservicios del sector bancario, donde aislar las reglas de negocio de los
adaptadores de infraestructura facilitaba su mantenimiento y su sustitución.
En VeriVoto, el dominio de cada servicio no depende de Spring ni de la base de
datos, lo que permite probar reglas electorales como el ciclo de vida del
proceso, el consumo del token de emisión o la verificación de la cadena de
auditoría mediante pruebas unitarias, sin levantar infraestructura; las pruebas
de arquitectura comprueban que esa dependencia no se introduce. De la misma
experiencia procede la decisión contraria respecto a la mensajería: conocer el
coste de operar Apache Kafka llevó a descartarlo en el MVP y a optar por REST
con un patrón outbox (D-06), suficiente para el volumen previsto.
